\chapter{Dual-Core DIP Example}
\label{chap:fpgasteel-dualcore-dip}

\section{About this project}
\label{sec:fpgasteel-dualcore-dip-about}

\textit{``Dual-Core DIP Example''} takes \nameref{chap:fpgasteel-dualcore-forwarding}'s camera-to-\gterm{vga}{VGA} pipeline and replaces its plain format conversion with a small computer-vision pipeline, \ccode{DualCoreDIP}: on every captured frame it looks for one white ball against a black background and, if it finds one that is not touching the image border, draws its outline and a centroid cross on the \gterm{vga}{VGA} output. It is the last example of this part, and the only one in which the two cores share an actual image-analysis workload rather than a mere format conversion.

It is also where this part measures what its last two optimizations are worth: the same algorithm is compiled three times, differing only in whether the \term{L1} instruction cache and \gterm{neon}{NEON} instructions are used. The three are \repopath{dualcore_dip_example}, \repopath{dualcore_dip_example_no_neon} and \repopath{dualcore_dip_example_no_neon_no_icache} (all under \repopath{sw/}); Paragraph~\ref{subsec:fpgasteel-dualcore-dip-build-variants} compares them, and everything else in this chapter holds for all three.

At a macroscopic level, the pipeline:

\begin{enumerate}
 \item converts the captured \gterm{rgb565}{RGB565} frame coming from the camera into two outputs, produced together in the same pass: an \term{RGB24} frame for \gterm{vga}{VGA} display, and a grayscale frame for analysis;
 \item thresholds it to a binary image (a pixel counts as ``object'' above a fixed brightness);
 \item cleans the binary image up morphologically, an opening (erosion followed by dilation) to remove isolated noise pixels without changing the blob's overall size;
 \item rejects the frame early if the cleaned blob touches the image border (it cannot be fully measured) or if too few pixels survived thresholding (nothing there);
 \item extracts the blob's outline, bounding box, and centroid, and overlays all three on the \term{RGB24} frame the camera pipeline already produced.
\end{enumerate}

\repopath{main0.cpp} reports the outcome on the \term{LED}s (\term{LED0} = object found, \term{LED1} = touching the border, \term{LED2} = empty) and the per-frame processing time on the hex display, exactly as \nameref{chap:fpgasteel-dualcore-forwarding} reports its own conversion time.

Every stage above runs split across both cores, using the exact same \ccode{Worker}/\ccode{Core1} dispatch mechanism documented in Paragraph~\ref{subsec:fpgasteel-dualcore-forwarding-dispatch}: nothing new had to be invented here, this project is simply that mechanism invoked several times instead of once. Section~\ref{sec:fpgasteel-dualcore-dip-operating-principles} covers the pipeline and its cache discipline; the rest of this chapter covers only the one genuinely new class, \ccode{DualCoreDIP}, since everything else (\ccode{Camera}, \ccode{Vga}, \ccode{CacheAwareFrame}/\ccode{CacheAwareFramePool}, \ccode{Core1}, \ccode{Worker}, \ccode{DualCoreJob}, \ccode{System}) is unchanged from \nameref{chap:fpgasteel-dualcore-forwarding}.

\section{Prerequisites}
\label{sec:fpgasteel-dualcore-dip-prerequisites}

Same toolchain, \term{SD card} and \term{Eclipse}/\gterm{openocd}{OpenOCD} debug setup as \nameref{chap:fpgasteel-dualcore-forwarding}, including the plain \repopath{boot.script.jtag_debug} script and the same two \term{CPU1}-related debug-configuration requirements.

\begin{warningbox}
\textbf{Hardware bug:} do not leave the camera in \term{RESET} or \term{POWER-DOWN} for extended periods.\footnotemark
\end{warningbox}
\footnotetext{See Part~\ref{part:appendices}, Section~\ref{sec:camera-adapter-reset-powerdown}.}

\section{FPGA Hardware Involved}
\label{sec:fpgasteel-dualcore-dip-hardware}

Unchanged: see Section~\ref{sec:fpgasteel-cache-forwarding-hardware}. Unlike \nameref{chap:fpgasteel-dualcore-bringup} and \nameref{chap:fpgasteel-dualcore-forwarding}, this project actually drives the \hw{fpga_led_pio}: after each processed frame, it lights a different \term{LED} depending on whether the ball was found, found touching the border, or not found at all, giving a live visual readout of the pipeline's outcome.

\section{Operating principles}
\label{sec:fpgasteel-dualcore-dip-operating-principles}

\subsection{The pipeline, stage by stage}
\label{subsec:fpgasteel-dualcore-dip-pipeline}

\ccode{DualCoreDIP::process()} runs several parallel stages and, when not convenient to parallelize, some single-core ones, to obtain the result introduced in Section~\ref{sec:fpgasteel-dualcore-dip-about}, in a fixed sequence, always in the same order, one synchronization point after another:

\begin{figure}[htbp]
 \centering
 \includesvg[width=\linewidth]{FPGAsteelDipPipeline}
 \caption{Every green stage is one \texttt{forkJoin()} (both cores), every blue stage runs on \term{CPU0} alone. The two early exits are what let an empty or unusable scene skip most of the pipeline instead of paying for all eleven stages every frame.}
 \label{fig:fpgasteel-dualcore-dip-pipeline}
\end{figure}

For each parallel step, inside the \ccode{process()} method, the code first prepares the two \ccode{Worker}s, \ccode{m_worker0} and \ccode{m_worker1}, by assigning to each one what function to run and what data to run it on (e.g., one half for each core). \ccode{process()} itself, like every other function involved, always runs on \term{CPU0}; it never runs on \term{CPU1}. The fork is \ccode{DualCoreDIP::forkJoin()} handing \ccode{m_worker1} to \term{CPU1} through \ccode{Core1::dispatch()}. It then immediately continues on \term{CPU0} with \ccode{m_worker0}'s own half of the work. Once \term{CPU0} finishes its own half, it starts polling \ccode{m_worker1}'s status until it reports it done: this is the join, and what makes \ccode{forkJoin()} a genuine synchronization point between one stage and the next.

\subsection{The three build variants}
\label{subsec:fpgasteel-dualcore-dip-build-variants}

As anticipated in Section~\ref{sec:fpgasteel-dualcore-dip-about}, the three sibling projects differ in whether they use the instruction cache and \gterm{neon}{NEON} instructions (see Table~\ref{tab:fpgasteel-dualcore-dip-build-variants}). \gterm{neon}{NEON} is \term{ARM}'s \term{SIMD} (Single Instruction, Multiple Data) instruction set; here, it speeds up loading pixel blocks into the \term{CPU}'s registers and processing them. During development, the compiler was found to auto-vectorize only sporadically, even when explicitly instructed to; for this reason, every stage in Figure~\ref{fig:fpgasteel-dualcore-dip-pipeline} except the two single-core ones (vii and xi) was rewritten by hand using \term{ARM} Intrinsics.

\begin{table}[htbp]
 \centering
 \begin{tabularx}{\textwidth}{@{} l X l @{}}
  \toprule
  \textbf{Project}\footnotemark & \textbf{Step implementations} & \textbf{L1 instruction cache} \\
  \midrule
  \repopath{_no_neon_no_icache} & plain scalar C++, same algorithm & disabled \\
  \lightmidrule
  \repopath{_no_neon} & plain scalar C++, same algorithm & enabled \\
  \lightmidrule
  (none) & \term{ARM} \gterm{neon}{NEON} intrinsics & enabled \\
  \bottomrule
 \end{tabularx}
 \caption{The three build variants of this project.}
 \label{tab:fpgasteel-dualcore-dip-build-variants}
\end{table}
\footnotetext{All three share the \texttt{dualcore\_dip\_example} prefix; only the distinguishing suffix is shown here.}

Since the board's seven-segment display reports each frame's processing time in milliseconds (see \repopath{main0.cpp}), comparing that figure across the three build variants is enough to work out how much faster one option is than another. Table~\ref{tab:fpgasteel-dualcore-dip-speedup} reports the resulting speedup: as shown, enabling the instruction cache dramatically cuts execution time (142 ms down to 47 ms), while the cores' optimized use through \gterm{neon}{NEON} instructions contributes a further reduction, smaller in absolute terms than the first, but still a significant 30 ms.

\begin{table}[htbp]
 \centering
 \begin{tabularx}{\textwidth}{@{} X l l l l @{}}
  \toprule
  \textbf{Project} & \textbf{L1 instruction cache} & \textbf{NEON} & \textbf{t [ms]} & \textbf{Speedup} \\
  \midrule
  \repopath{_no_neon_no_icache} & disabled & unused & 142 & --- (baseline) \\
  \lightmidrule
  \repopath{_no_neon} & enabled & unused & 47 & 3.02$\times$ \\
  \lightmidrule
  (none) & enabled & explicit & 17 & 8.35$\times$ \\
  \bottomrule
 \end{tabularx}
 \caption{Measured per-frame processing time and speedup across the three variants. Isolating \term{NEON}'s own contribution yields a further 2.76$\times$ speedup on top of the instruction cache's.}
 \label{tab:fpgasteel-dualcore-dip-speedup}
\end{table}


\section{How it works}
\label{sec:fpgasteel-dualcore-dip-how-it-works}

The following snippet is taken from \repopath{main0.cpp}:

\begin{lstlisting}[language=C++]
sys->prepareCacheSubsystem();
Core1 & core1 = Core1::getInstance();
core1.start();
while (core1.getStatus() != Core1::Status::Ready
 && core1.getStatus() != Core1::Status::Error) {}
CacheAwareFramePool camPool(CAMERA_POOL_SIZE, FRAME_WIDTH, FRAME_HEIGHT,
 CAMERA_COLOR_SCHEME);
CacheAwareFramePool vgaPool(VGA_POOL_SIZE, FRAME_WIDTH, FRAME_HEIGHT,
 Frame::ColorScheme::RGB24);
vga->init(vgaPool);
pclkResetCtrl->deassertReset();
sys->delay(1000);
camera->init(camPool);
hex->showDashes();
DualCoreDIP dip(FRAME_WIDTH, FRAME_HEIGHT);
while (true) {
 if (camPool.available() > 0) {
  CacheAwareFrame * camFrame = camPool.pop(true);
  CacheAwareFrame * vgaFrame = static_cast<CacheAwareFrame *>(vgaPool.borrow());
  uint64_t t0 = sys->millis();
  DualCoreDIP::Result res = dip.process(*camFrame, *vgaFrame);
  uint64_t t1 = sys->millis();
  camPool.giveBack(camFrame);
  vgaPool.push(vgaFrame, true);
  LED_PIO = static_cast<uint8_t>(res);
  // LED0=OK, LED1=ON_BORDER, LED2=EMPTY
  hex->showNumber(static_cast<uint32_t>(t1 - t0));
 }
}
\end{lstlisting}

Structurally this is Section~\ref{sec:fpgasteel-dualcore-forwarding-how-it-works} with \ccode{dip.process(*camFrame, *vgaFrame)} standing in for \ccode{converter.convert(*camFrame, *vgaFrame)}, plus the one line writing the result code to \hw{LED_PIO}. The \gterm{exception}{exception} list, pool handling, and cache-flush placement around \ccode{pop(true)}/\ccode{push(vgaFrame, true)} are all identical.

\section{Debug tooling}
\label{sec:fpgasteel-dualcore-dip-debug-tooling}

\repopath{debugTools/view_frame.py}, same tool as Section~\ref{sec:fpgasteel-cache-forwarding-debug-tooling}. This project additionally ships \repopath{debugTools/linux_dualcore_emulator}, a small host-side (\term{Linux}, x86/\term{ARM}) build of \ccode{DualCoreDIP}'s pipeline logic against stub \gterm{hwlib}{hwlib} headers, used to test and tune the vision algorithm itself (thresholds, morphology, early-exit conditions) against saved frames on a development machine, without needing the board at all; it is not part of the bare-metal build and is not covered further here.